%% kapa: exam=1 round=28 year=2017 period=2 form=A subject=appraisal_law subject_ko=감정평가관계법규 seq=15 number=15
%% source: Q-Net 감정평가사 기출문제 (cst00309) artlSeq=5199603 file=extracted/1008/5199603/2017년 제28회 감정평가사 1차시험 문제지/2교시 A형.pdf
%% conversion: kordoc 4.24.51 -> tools/gen1.py (auto); review=proofread-sonnet; flags=none
\begin{question}{15}
감정평가 및 감정평가사에 관한 법령상 감정평가사의 권리와 의무에 관한 설명으로 옳지 않은 것은?
\begin{choices}
  \item 감정평가사는 2명 이상의 감정평가사로 구성된 합동사무소를 설치할 수 있다.
  \item 감정평가사가 감정평가업을 하려는 경우에는 감정평가사무소 개설에 대하여 국토교통부장관의 인가를 받아야 한다.
  \item 감정평가사는 감정평가업을 하기 위하여 1개의 사무소만을 설치할 수 있다.
  \item 감정평가업자는 토지등의 매매업을 직접 하여서는 아니 된다.
  \item 감정평가업자가 손해배상책임을 보장하기 위하여 보증보험에 가입하는 경우 보험 가입 금액은 감정평가사 1인당 1억원 이상으로 한다.
\end{choices}
\end{question}
\answer{2}
